States are the occupation states $e_{S}$, $S \subseteq \{1, \dots, 5\}$, with slot number $j = |S \cap \{4, 5\}|$ and pair number $k = |S \cap \{1, 2, 3\}|$, $B - L = 1 - \tfrac23k$, $Y = \tfrac12(j - 1) + \tfrac12(B - L)$, $T_{3} = \tfrac12(N_{4} - N_{5})$, as in \claimref{synthesis.abelian-charges-of-the-sixteen-and-the-terms-that-break-them} under Hypothesis (S) (\claimref{standard-model.hyp-s-slot-statistics}). In this labelling the neutral member of $L$ is $e_{\{1,2,3,4\}}$ and the electrically neutral slot operators are $v_{5}$, $w_{5}$; which rung is called neutral is a convention.

\emph{(i).} $B - L$ is a function of the pair number, so an operator that shifts $k$ by a fixed amount $d$ is a weight vector of weight $-\tfrac23d$. A pair-pair bilinear shifts $k$ by $0$ ($v_{a}w_{b}$, $w_{a}v_{b}$) or $\pm2$; a slot-pair bilinear by $\pm1$; a slot-slot bilinear and a slot vector by $0$; $T = v_{1}\cdots v_{5}$ by $+3$. This is the membership table of \claimref{synthesis.baryon-number-centraliser-in-the-five-mode-algebra}(ii) to (iv), with $[B - L, T] = -2T$ (\claimref{standard-model.majorana-element}); the check recounts the weights on all $32$ states: ten, twelve and twelve, three and three. A gauge term couples a generator to its gauge field, which transforms in the adjoint and so carries the opposite weight: the term is invariant under the rotation generated by $B - L$. A rung vector has weight $0$ whether it is a field or a vacuum direction. $A_{2}(B - L) = A_{3}(B - L) = 0$ and the vanishing trace on the sixteen are clause (v) of the first cited claim, with the coefficients of Fact \claimref{standard-model.imported-anomaly-coefficients}. The traces over species: $\tr(B - L)$ is $6 \cdot \tfrac13 - 3 \cdot \tfrac13 - 3 \cdot \tfrac13 - 2 + 1 = -1$ on the fifteen and $0$ with the singlet's $+1$; $\tr((B - L)^{3}) = \tfrac{6}{27} - \tfrac{3}{27} - \tfrac{3}{27} - 2 + 1 = -1$ and $0$; $\tr(Y^{2}(B - L)) = \tfrac1{18} - \tfrac49 - \tfrac19 - \tfrac12 + 1 = 0$ and $\tr(Y(B - L)^{2}) = \tfrac19 - \tfrac29 + \tfrac19 - 1 + 1 = 0$, on the fifteen and on the sixteen alike since the singlet has $Y = 0$. The trace of a charge is its pairing with the gravitational four-form in the same sense in which $A_{2}$ and $A_{3}$ are its pairings with the weak and colour four-forms (the mixed gravitational coefficient of the cited Fact, read for $B - L$ in place of $Y$). $T$ multiplies no field in the stored chain: it is the direction of a vacuum value, so its weight $-2$ is not compensated.

\emph{(ii).} A bilinear form on the sixteen is a combination of the $256$ forms $\beta_{S,T}(\psi, \varphi) = \psi_{S}\varphi_{T}$, and $\beta_{S,T}$ has weight $x_{S} + x_{T}$ under a diagonal charge $x$. Invariance under the electric charge and the two diagonal colour generators keeps the pairs of zero total weight; invariance under the six colour root vectors is then a linear system on those, whose kernel the check computes: dimension $10$, each kernel vector supported on one pair of species. By hand: a nonzero $B - L$ weight needs two leptons, two quarks of the same triality, or a quark and a lepton; $3 \otimes 3$, $\bar{3} \otimes \bar{3}$ and $3 \otimes 1$ contain no colour singlet, so only lepton pairs remain, and among them the neutral ones are $(\nu, \nu)$, $(\nu, \nu^{c})$, $(\nu^{c}, \nu^{c})$, $(e, e^{c})$ with weights $-2, 0, +2, 0$. The form on $(\nu^{c}, \nu^{c})$ is $B(\psi, T\varphi)$: $Te_{\emptyset} = e_{\{1,\dots,5\}}$, the complement of the empty state, and $T$ annihilates every other even state (\claimref{standard-model.invariant-pairing-five-modes}). The form on $(\nu, \nu)$ has $T_{3} = 1$: it is a component of a weak triplet and is not invariant under $\su(2)_{L}$.

\emph{(iii).} By (i) each interaction of the list conserves $B - L$ when the gauge fields are assigned their weights. The anomalous weak process of (A) changes a charge $X$ in proportion to $A_{2}(X)$ (Fact \claimref{synthesis.imported-anomalous-weak-process-changes-b-plus-l-and-keeps-b-minus-l}, clause (a) and its reading), and the colour process, which is literature and not stored, in proportion to $A_{3}(X)$; both coefficients vanish for $B - L$. The gravitational pairing does not vanish on the fifteen, by (i); it is closed by (Gr). That the gravitational four-form vanishes on a homogeneous isotropic background is the statement that it is built from the Weyl part of the curvature, which vanishes there; this is standard and argued, not derived in the store, and the size of the channel on a perturbed background is not computed. The remaining stored term is the Majorana element. Its quanta, the heavy neutral singlets, are absent by (R) (\claimref{neutrino.r-descent-reheats-below-lower-heavy}); its effect at second order on the fifteen is excluded by (W); its non-thermal effect by (I), which holds under \claimref{neutrino.real-breaking-assumption}. By (ii) no further bilinear term with a fixed coefficient exists that respects colour and electric charge and changes $B - L$, except the triplet form, which the stored chain does not contain. So the total $B - L$ is constant after the descent; that this total is the registered datum is the naming (N) (\claimref{synthesis.matter-asymmetry-registered-free-input}). The inference from the listed terms to the dynamics of the plasma is the usual one from a symmetry of the action to a conserved charge and is not derived in the store.

\emph{(iv).} The conditions are linear. Colour: the charge is constant on colour multiplets. The weak root vector $v_{4}w_{5}$ carries a charge $h_{W}$: $x_{T} - x_{S} = h_{W}$ whenever it maps $e_{S}$ to $e_{T}$. A Dirac pairing $B(\psi, a\varphi)$ with a slot operator $a$ of charge $h_{a}$ is conserved if $x_{S} + x_{T} + h_{a} = 0$ on each pair it couples; one doublet and its conjugate means $h(w_{s}) = -h(v_{s})$. The weak anomalous process involves, per family, one member of each weak doublet of each colour and of the lepton doublet with total $T_{3} = 0$; given the root-vector condition, its condition is that the sum of the charge over the eight doublet states vanishes, which is $A_{2}(X) = 0$ for charges commuting with the weak algebra. The check computes the kernels: dimensions $4$ and $3$, and in the broken phase, where $h(v_{5}) = h(w_{5}) = 0$, $3$ and $2$, with the stated bases. That independent charges $h_{v}$, $h_{w}$ would add exactly one charge, with $A_{3} = -2(h_{v} + h_{w}) \neq 0$, is clause (ii) of the first cited claim; with the colour anomalous process in equilibrium it is excluded, so $h_{w} = -h_{v}$ for every doublet that couples to the quarks.

\emph{(v).} In a grand canonical state with chemical potentials $\mu$ for the conserved charges, a massless fermion state of charge $x$ has excess density $\tfrac{T^{2}}{6}x\mu$ to first order and a complex boson state twice that; the two numbers are $\tfrac{1}{\pi^{2}}\int_{0}^{\infty}p^{2}\,\partial_{\mu}\bigl(f(p - \mu) - f(p + \mu)\bigr)dp$ at $\mu = 0$ for the Fermi and Bose distributions, with $\int_{0}^{\infty}t^{2}e^{t}(e^{t} \pm 1)^{-2}dt = \tfrac{\pi^{2}}{6}, \tfrac{\pi^{2}}{3}$. Hence $n_{X} = \tfrac{T^{2}}{6}\chi(X, \mu)$ with $\mu$ read as a charge, an element of the space $V$ of (iv); this is the thermal input (T) and is standard. $\chi$ is positive definite. Neutrality is $\chi(g, \mu) = 0$ for the gauge charges $g$ of the phase, which leaves a line in $V$: $V = \mathrm{span}\{T_{3}, Y, B - L\}$ with two conditions, or $\mathrm{span}\{Q_{\mathrm{em}}, B - L\}$ with one. In the symmetric phase $\chi(T_{3}, Y) = \chi(T_{3}, B - L) = \chi(T_{3}, n_{B}) = 0$, since $Y$, $B - L$, $n_{B}$ are constant on weak doublets, so $T_{3}$ drops out. With one gauge charge $g$, $\mu \propto \chi(g, g)(B - L) - \chi(g, B - L)g$, and the displayed ratio follows. The traces over the fifteen: $\tr(Y^{2}) = \tfrac16 + \tfrac43 + \tfrac13 + \tfrac12 + 1 = \tfrac{10}{3}$; $\tr(Y(B - L)) = \tfrac13 + \tfrac23 - \tfrac13 + 1 + 1 = \tfrac83$; $\tr((B - L)^{2}) = \tfrac{12}{9} + 3 = \tfrac{13}{3}$; $\tr(n_{B}(B - L)) = \tfrac43$; $\tr(n_{B}Y) = \tfrac23$; $\tr(Q_{\mathrm{em}}^{2}) = \tfrac{16}{3}$; and $T_{3}$ is trace-orthogonal to $B - L$ and $n_{B}$. A doublet adds $2 \cdot 2 \cdot \tfrac14 = 1$ to $\chi(Y, Y)$; in the broken phase the weak boson of charge one, counted with two helicities, adds $4$ to $\chi(Q_{\mathrm{em}}, Q_{\mathrm{em}})$ and each charged scalar $2$ (for one doublet the charged scalar is the third polarisation of the weak boson). This count, massless and with the charged component of every doublet a state of the plasma, is the convention of the literature's closed form (Fact \claimref{synthesis.imported-equilibrium-baryon-number-as-a-fraction-of-b-minus-l}, clauses (d), (e)); with three polarisations and a charged scalar in addition the boson term would be $8$ at $N_{H} = 1$ and the ratio $\tfrac{10}{31}$. Then
\[
c_{\mathrm{sym}} = \frac{\tfrac43(\tfrac{10}{3}N_{f} + N_{H}) - \tfrac{16}{9}N_{f}}{\tfrac{13}{3}(\tfrac{10}{3}N_{f} + N_{H}) - \tfrac{64}{9}N_{f}} = \frac{24N_{f} + 12N_{H}}{66N_{f} + 39N_{H}}, \qquad c_{\mathrm{br}} = \frac{\tfrac43(\tfrac{16}{3}N_{f} + 4 + 2N_{H}) - \tfrac{16}{9}N_{f}}{\tfrac{13}{3}(\tfrac{16}{3}N_{f} + 4 + 2N_{H}) - \tfrac{64}{9}N_{f}} = \frac{48N_{f} + 24(N_{H} + 2)}{144N_{f} + 78(N_{H} + 2)},
\]
which reduce to the stated forms. The check evaluates the Gram ratio from the states on the grid $N_{f} = 1, \dots, 6$, $N_{H} = 0, \dots, 6$, and compares $c_{\mathrm{sym}}$ with the kernel of the five conditions used by the check of Computation \claimref{dark-tower-matter.mirrors-composition}.

\emph{(vi).} Put $\Delta_{f} = \tfrac13n_{B} - L_{f}$, so $\sum_{f}\Delta_{f} = B - L$, and $\mu = aY + \sum_{f}b_{f}\Delta_{f}$. $\chi(Y, \Delta_{f})$ and $\chi(n_{B}, \Delta_{f})$ do not depend on $f$, and $\sum_{f}\chi(\Delta_{f}, \Delta_{f'})$ does not depend on $f'$. So neutrality fixes $a$ in terms of $\sum_{f}b_{f}$, and $n_{B}$ and $n_{B-L} = \sum_{f}n_{\Delta_{f}}$ are both multiples of $\sum_{f}b_{f}$: their ratio is $c$.

\emph{(vii).} In the covariant coupling the quark mass matrices are $M_{u} = \alpha Y + \bar{\beta}\tilde{Y}$, $M_{d} = \beta Y + \bar{\alpha}\tilde{Y}$ (\claimref{synthesis.rung-field-coupling-to-its-conjugate-is-the-second-family-matrix}): the doublets $\alpha$ and $\bar{\beta}$ enter the same pairings, so a conserved charge assigns them the same value, and each present doublet adds one to $N_{H}$. The lower bound of about $2.5$ TeV on the second doublet with a level is the conclusion of \claimref{synthesis.covariant-rung-coupling-gives-the-normal-pair-tree-level-b-meson-mixing-that-excludes-the-transmutation-branch}; a species of mass $m$ contributes to $\chi$ with a weight that falls exponentially for $m \gg T$. A bound: for a boson of mass $xT$ the weight relative to a massless one is $\tfrac{3}{\pi^{2}}\int_{0}^{\infty}t^{2}e^{\varepsilon}(e^{\varepsilon} - 1)^{-2}dt$ with $\varepsilon = \sqrt{t^{2} + x^{2}}$; since $e^{\varepsilon}(e^{\varepsilon} - 1)^{-2} \le e^{-\varepsilon}(1 - e^{-x})^{-2}$ and $5\varepsilon \ge 3t + 4x$, it is at most $\tfrac{3}{\pi^{2}}(1 - e^{-x})^{-2}\,2(\tfrac53)^{3}e^{-4x/5} \le \tfrac{250}{81}(1 - 2^{-18})^{-2}e^{-4x/5}$ for $x \ge 18$, using $\pi^{2} > 9$ and $e^{-18} < 2^{-18}$. With a second doublet at weight $\epsilon$ the Gram ratio is $(28 + 4\epsilon)/(79 + 13\epsilon) = \tfrac{28}{79} - 48\epsilon/(79(79 + 13\epsilon))$. The stored bound of about $2500$ GeV and a freeze-out at most $134$ GeV give $x \ge 18$, a weight below $2 \times 10^{-6}$ and a shift of $c$ below $2 \times 10^{-8}$; the check computes these from the Gram ratio. The inputs to this are a stored bound stated as approximate and a literature temperature, which is why the count is argued. That baryon number follows the equilibrium value until the weak process freezes out, and that it freezes out at $T_{*} = 131.7 \pm 2.3$ GeV, below the crossover at $159 \pm 1$ GeV, is imported (Fact \claimref{synthesis.imported-equilibrium-baryon-number-as-a-fraction-of-b-minus-l}, clause (f); Fact \claimref{synthesis.imported-lattice-crossover-and-freeze-out-of-the-anomalous-weak-process}). The interpolation, clause (a) of the first of these, is a mediant: $(28a + 72b)/(79a + 222b)$ with $a, b \ge 0$ lies between $\tfrac{28}{79}$ and $\tfrac{72}{222} = \tfrac{12}{37}$. The bounds on the datum are $\tfrac{79}{28}$ and $\tfrac{37}{12}$ times $8.7 \times 10^{-11}$.
